\begin{titlepage}
\thispagestyle{empty}
\begin{tikzpicture}[remember picture, overlay]
  \fill[inkdark] (current page.north west) rectangle (current page.south east);
  \fill[edugreen!25] ([xshift=-0.6cm,yshift=-2.0cm]current page.north east) circle (6.3cm);
  \fill[white!8] ([xshift=1.2cm,yshift=2.0cm]current page.south west) circle (5.0cm);

  \draw[honeygold!85, line width=1.4pt]
    ([xshift=1.1cm,yshift=-1.1cm]current page.north west)
    rectangle
    ([xshift=-1.1cm,yshift=1.1cm]current page.south east);

  \node[
    anchor=north west,
    fill=honeygold,
    text=inkdark,
    font=\bfseries\large,
    rounded corners=4pt,
    inner xsep=10pt,
    inner ysep=6pt
  ] at ([xshift=1.6cm,yshift=-1.5cm]current page.north west) {QUYỂN XVIII};

  \node[
    anchor=north,
    text=white,
    font=\large\itshape
  ] at ([yshift=-1.7cm]current page.north) {Truyện Tích Cảm Hứng Song Ngữ};

  \node[
    anchor=west,
    text=white,
    font=\fontsize{22}{26}\selectfont\bfseries
  ] at ([xshift=1.8cm,yshift=4.2cm]current page.center) {NHỮNG LỜI};

  \node[
    anchor=west,
    text=honeygold,
    font=\fontsize{22}{26}\selectfont\bfseries
  ] at ([xshift=1.8cm,yshift=2.3cm]current page.center) {THẦY NÓI};

  \node[
    anchor=west,
    text=white,
    font=\fontsize{22}{26}\selectfont\bfseries
  ] at ([xshift=1.8cm,yshift=0.4cm]current page.center) {TRONG LỚP HỌC};

  \draw[honeygold, line width=2pt]
    ([xshift=1.8cm,yshift=-1.0cm]current page.center) --
    ([xshift=8.8cm,yshift=-1.0cm]current page.center);

  \node[
    anchor=west,
    text=white,
    font=\large
  ] at ([xshift=1.8cm,yshift=-2.1cm]current page.center) {50 câu nói ngắn, sâu — mỗi câu một câu chuyện nhỏ — để mở đầu bài, kết tiết, hoặc tặng học sinh};

  \node[
    anchor=north west,
    fill=white,
    text=black,
    rounded corners=8pt,
    text width=11.7cm,
    inner sep=14pt,
    align=left,
    font=\normalsize
  ] at ([xshift=1.9cm,yshift=-14.9cm]current page.north west) {%
    \textbf{Tinh thần quyển này:} gom 50 câu nói rất ``đắt'' của thầy giáo trong lớp học — ngắn, sâu, học sinh nghe xong nhớ lâu. Mỗi câu đi kèm một câu chuyện nhỏ để gợi suy nghĩ.\\[6pt]
    \textbf{Cách dùng:} mở đầu bài học, kết thúc tiết học, viết lên bảng, chèn vào câu chuyện giáo dục, hoặc tặng học sinh cuối năm.\\[6pt]
    \textbf{Điểm tựa:} không giảng đạo lý dài dòng; mỗi trang một câu, một chuyện, một bài học giữ lại.
  };

  \node[
    anchor=south,
    text=honeygold,
    font=\large\itshape,
    align=center
  ] at ([yshift=2.1cm]current page.south) {%
    ``Học không phải để hơn người khác,\
    mà để hơn chính mình hôm qua.''%
  };

  \node[
    anchor=south,
    text=white!85,
    font=\normalsize,
    align=center
  ] at ([yshift=1.1cm]current page.south) {Dành cho giáo viên và học sinh — những lời thầy nói trong lớp học};
\end{tikzpicture}
\end{titlepage}
\clearpage
